\section*{अध्याय \ref{ch:g12:proton-nmr} --- प्रोटॉन NMR}

\begin{solution}{exo:g12:proton-nmr:1}
मेथेन: 1। एथेन: 1 (दोनों \ce{CH3} एक जैसे हैं)। प्रोपेन: 2 (दो
\ce{CH3} एक जैसे, एक \ce{CH2})। मेथेनॉल: 2 (\ce{CH3} और \ce{OH})।
मेथॉक्सीमेथेन: 1।
\end{solution}

\begin{solution}{exo:g12:proton-nmr:2}
2 पड़ोसी: एक त्रिक, 1:2:1। 3 पड़ोसी: एक चतुष्क, 1:3:3:1।
0 पड़ोसी: एक एकक।
\end{solution}

\begin{solution}{exo:g12:proton-nmr:3}
8 \omterm{def:g9:inside-the-atom:proton}{प्रोटॉनों} के लिए कुल \qty{40}{mm}: प्रति \omterm{def:g9:inside-the-atom:proton}{प्रोटॉन} \qty{5}{mm}। संकेत
3, 2 और 3 \omterm{def:g9:inside-the-atom:proton}{प्रोटॉनों} को दर्शाते हैं।
\end{solution}

\begin{solution}{exo:g12:proton-nmr:4}
\omterm{def:g11:functional-groups:carbonyl}{ऐल्डिहाइड}: 9.7 से \qty{10.0}{ppm}। श्रृंखला का \ce{CH3}: 0.7 से
\qty{1.3}{ppm}।
\end{solution}

\begin{solution}{exo:g12:proton-nmr:5}
उसके बारह \omterm{def:g12:proton-nmr:equivalent}{तुल्य प्रोटॉन} एक अकेली, तीक्ष्ण, प्रबल रेखा देते हैं, जिसे
शून्य पर रखना आसान है। ड्यूटेरीकृत \omterm{def:g6:solutions-and-solubility:solute}{विलायक} अपना कोई \omterm{def:g9:inside-the-atom:proton}{प्रोटॉन} संकेत नहीं
देता, जो अन्यथा नमूने के कुछ मिलीग्राम के संकेतों को ढक लेता।
\end{solution}

\begin{solution}{exo:g12:proton-nmr:6}
दो संकेत। क्लोरीन के बगल का \ce{CH2}: 2 \omterm{def:g9:inside-the-atom:proton}{प्रोटॉन}, एक चतुष्क (3
पड़ोसी), 2.5--\qty{4.0}{ppm} के परास में। \ce{CH3}: 3 \omterm{def:g9:inside-the-atom:proton}{प्रोटॉन}, एक
त्रिक (2 पड़ोसी), \ce{CH3} के सामान्य परास के निकट, थोड़ा ऊपर,
क्योंकि क्लोरीन दो \omterm{def:g7:atoms-and-molecules:atom}{परमाणु} दूर है।
\end{solution}

\begin{solution}{exo:g12:proton-nmr:7}
तीन संकेत: दोनों \ce{CH3}, 6 \omterm{def:g12:proton-nmr:equivalent}{तुल्य प्रोटॉन}, एक द्विक (1
पड़ोसी, \ce{CH}); \ce{CH}, 1 \omterm{def:g9:inside-the-atom:proton}{प्रोटॉन}, 6 पड़ोसियों द्वारा
7 रेखाओं में विभक्त, \ce{H-C-O} परास (3.3--\qty{4.5}{ppm}) में;
\ce{OH}, 1 \omterm{def:g9:inside-the-atom:proton}{प्रोटॉन}, एक एकक।
\end{solution}

\begin{solution}{exo:g12:proton-nmr:8}
(a) प्रोपेनोन; (b) एथिल एथेनोएट; (c) प्रोपेनोइक अम्ल।
\end{solution}

\begin{solution}{exo:g12:proton-nmr:9}
\ce{CH2} \omterm{def:g11:functional-groups:carboxylic-acid}{एस्टर} के ऑक्सीजन \omterm{def:g7:atoms-and-molecules:atom}{परमाणु} से आबंधित है, जो
विद्युत-ऋणात्मक है और उसके \omterm{def:g9:inside-the-atom:proton}{प्रोटॉनों} का परिरक्षण घटाता है: \ce{H-C-O} परास,
3.3--\qty{4.5}{ppm}। \ce{CH3} एक कार्बन और दूर है और
श्रृंखला के परास में ही रहता है।
\end{solution}

\begin{solution}{exo:g12:proton-nmr:10}
प्रोपेनोन: एक \ce{C=O} (\omterm{def:g11:functional-groups:carbonyl}{कीटोन}, \qty{1715}{cm^{-1}}) और छह \omterm{def:g12:proton-nmr:equivalent}{तुल्य
प्रोटॉन}, एक एकक। प्रोपेनैल तीन संकेत देता, जिनमें से एक,
\ce{CH3-CH2-CHO} का ऐल्डिहाइडी \omterm{def:g9:inside-the-atom:proton}{प्रोटॉन}, लगभग 9.7--\qty{10.0}{ppm} पर होता।
\end{solution}

\begin{solution}{exo:g12:proton-nmr:11}
\ce{OH} \omterm{def:g9:inside-the-atom:proton}{प्रोटॉन} का \omterm{def:g7:atoms-and-molecules:molecule}{अणुओं} के बीच तेज़ी से आदान-प्रदान होता है: वह
एकक देता है और अपने पड़ोसियों को विभक्त नहीं करता। इसलिए \ce{CH2}
केवल \ce{CH3} के 3 \omterm{def:g9:inside-the-atom:proton}{प्रोटॉनों} द्वारा विभक्त होता है: एक चतुष्क।
\end{solution}

\begin{solution}{exo:g12:proton-nmr:12}
ब्यूटेनोइक अम्ल \ce{CH3-CH2-CH2-COOH}: 4 संकेत, एक त्रिक (\ce{CH3}), एक
6-रेखा संकेत (बीच का \ce{CH2}, 5 पड़ोसी), एक त्रिक
(\ce{CH2-C=O}), लगभग 11--\qty{12}{ppm} पर एक एकक। एथिल एथेनोएट: एक
एकक, एक चतुष्क, एक त्रिक। मेथिल प्रोपेनोएट \ce{CH3-CH2-CO-O-CH3}:
एक त्रिक, एक चतुष्क (\ce{CH2-C=O}, 2.0--\qty{2.4}{ppm}), एक एकक
(\ce{O-CH3})। प्रोपिल मेथेनोएट \ce{H-CO-O-CH2-CH2-CH3}: 4 संकेत,
बहुत बाईं ओर एक एकक (\ce{C=O} वाले कार्बन पर H), एक त्रिक
(\ce{O-CH2}), एक 6-रेखा संकेत, एक त्रिक (\ce{CH3})।
\end{solution}

\begin{solution}{exo:g12:proton-nmr:13}
\ce{OH} संकेत: \ce{O-H} \omterm{def:g9:inside-the-atom:proton}{प्रोटॉन} भारी जल के ड्यूटीरियम से
आदान-प्रदान करते हैं, और \ce{O-D} कोई \omterm{def:g9:inside-the-atom:proton}{प्रोटॉन} संकेत नहीं देता। जो संकेत
\ce{D2O} के साथ हिलाने पर गायब हो जाए, वह O या N पर स्थित \omterm{def:g9:inside-the-atom:proton}{प्रोटॉन} का है।
\end{solution}

\begin{solution}{exo:g12:proton-nmr:14}
एथेनॉल लगभग \qty{3.69}{ppm} पर एक चतुष्क, लगभग
\qty{1.23}{ppm} पर एक त्रिक और एक \ce{OH} एकक जोड़ता है। \omterm{def:g11:functional-groups:carboxylic-acid}{एस्टर} का एकक: 3 \omterm{def:g9:inside-the-atom:proton}{प्रोटॉनों}
के लिए \qty{30}{mm}, \omterm{def:g11:functional-groups:carboxylic-acid}{एस्टर} के प्रति \omterm{def:g9:inside-the-atom:proton}{प्रोटॉन} \qty{10}{mm}। एथेनॉल का त्रिक:
3 \omterm{def:g9:inside-the-atom:proton}{प्रोटॉनों} के लिए \qty{3}{mm}, एथेनॉल के प्रति \omterm{def:g9:inside-the-atom:proton}{प्रोटॉन} \qty{1}{mm}। चूँकि हर
\omterm{def:g7:atoms-and-molecules:molecule}{अणु} इन संकेतों में एक \ce{CH3} का योगदान देता है, मात्राएँ
$1/10$ के अनुपात में हैं: \omterm{def:g11:functional-groups:carboxylic-acid}{एस्टर} के दस \omterm{def:g7:atoms-and-molecules:molecule}{अणुओं} पर एथेनॉल का लगभग एक \omterm{def:g7:atoms-and-molecules:molecule}{अणु}।
\end{solution}

\begin{solution}{exo:g12:proton-nmr:15}
चार संकेत: दोनों \ce{CH3} (6 \omterm{def:g9:inside-the-atom:proton}{प्रोटॉन}) एक द्विक; \ce{CH} (1
\omterm{def:g9:inside-the-atom:proton}{प्रोटॉन}), 6 + 2 = 8 पड़ोसियों द्वारा 9 रेखाओं में विभक्त; \ce{CH2} (2
\omterm{def:g9:inside-the-atom:proton}{प्रोटॉन}), \ce{O} के बगल में, एक द्विक, \ce{H-C-O} परास में;
\ce{OH} (1 \omterm{def:g9:inside-the-atom:proton}{प्रोटॉन}) एक एकक। \ce{CH} संकेत में सबसे अधिक रेखाएँ हैं।
\end{solution}

\begin{solution}{pb:g12:proton-nmr:1}
\textbf{1.} चार कार्बनों वाली \omterm{def:g11:organic-skeletons:alkane}{ऐल्केन} \ce{C4H10} है; एक \ce{C=O}
दो हाइड्रोजन घटाता है और दोनों ऑक्सीजन कुछ नहीं जोड़तीं: \ce{C4H8O2},
चार कार्बनों वाले अम्लों और \omterm{def:g11:functional-groups:carboxylic-acid}{एस्टरों} का सूत्र।

\textbf{2.} ब्यूटेनोइक अम्ल \ce{CH3-CH2-CH2-COOH}; एथिल एथेनोएट
\ce{CH3-CO-O-CH2-CH3}; मेथिल प्रोपेनोएट \ce{CH3-CH2-CO-O-CH3}; प्रोपिल
मेथेनोएट \ce{H-CO-O-CH2-CH2-CH3}।

\textbf{3.} तीनों \omterm{def:g11:functional-groups:carboxylic-acid}{एस्टर} \omterm{def:g11:functional-groups:carboxylic-acid}{एस्टर} समूह साझा करते हैं; चारों में एक
\ce{C=O} है।

\textbf{4.} एक \ce{C=O}, \omterm{def:g11:functional-groups:carboxylic-acid}{एस्टर} परास में (लगभग \qty{1735}{cm^{-1}})।

\textbf{5.} 2500 से \qty{3300}{cm^{-1}} तक बहुत चौड़ा \ce{O-H} पट्ट:
अनुपस्थित, इसलिए ब्यूटेनोइक अम्ल नहीं।

\textbf{6.} नहीं: तीनों \omterm{def:g11:functional-groups:carboxylic-acid}{एस्टरों} में वही समूह हैं।

\textbf{7.} तीन।

\textbf{8.} एक एथिल समूह \ce{CH3-CH2-}: \ce{CH2} 3
\omterm{def:g9:inside-the-atom:proton}{प्रोटॉनों} द्वारा विभक्त (चतुष्क), \ce{CH3} 2 द्वारा (त्रिक)।

\textbf{9.} उसके \omterm{def:g9:inside-the-atom:proton}{प्रोटॉनों} के पड़ोसी \omterm{def:g7:atoms-and-molecules:atom}{परमाणुओं} पर कोई \omterm{def:g9:inside-the-atom:proton}{प्रोटॉन} नहीं है: एक
\ce{CH3} जो \ce{C=O} वाले कार्बन से या \omterm{def:g11:functional-groups:carboxylic-acid}{एस्टर} के ऑक्सीजन से आबंधित है।

\textbf{10.} \ce{CH2} ऑक्सीजन से आबंधित है (\ce{H-C-O} परास)।

\textbf{11.} एक त्रिक (\ce{CH3}), लगभग 2.0--\qty{2.4}{ppm} पर एक चतुष्क
(\ce{CH2}, \ce{C=O} के बगल में) और \ce{H-C-O} परास में एक एकक
(\ce{O-CH3})। चतुष्क \qty{4.12}{ppm} से बहुत नीचे होता और
एकक \qty{2.04}{ppm} से बहुत ऊपर: यह अज्ञात यौगिक नहीं।

\textbf{12.} चार संकेत (बहुत बाईं ओर एक एकक, एक त्रिक, एक 6-रेखा
संकेत, एक त्रिक): यह अज्ञात यौगिक नहीं, जिसके तीन संकेत हैं।

\textbf{13.} एथिल एथेनोएट, \ce{CH3-CO-O-CH2-CH3}: एकक
\qty{2.04}{ppm} = \ce{CH3-C=O}; चतुष्क \qty{4.12}{ppm} = \ce{O-CH2};
त्रिक \qty{1.26}{ppm} = सिरे का \ce{CH3}।

\textbf{14.} एकक 3, चतुष्क 2, त्रिक 3।

\textbf{15.} प्रति \omterm{def:g9:inside-the-atom:proton}{प्रोटॉन} $72 / 8 = \qty{9}{mm}$।

\textbf{16.} एकक $3 \times 9 = \qty{27}{mm}$; त्रिक
\qty{27}{mm}।

\textbf{17.} $18 + 27 + 27 = \qty{72}{mm}$।

\textbf{18.} हाँ: छोटे \omterm{def:g11:functional-groups:carboxylic-acid}{एस्टर} अपनी फलों जैसी गंधों के लिए जाने जाते हैं।

\textbf{19.} चतुष्क की सीढ़ी \qty{72}{mm} में से \qty{18}{mm} है।
\end{solution}
